\documentclass[12pt]{article}
\usepackage[utf8]{inputenc}
\usepackage[T1]{fontenc}
\usepackage[margin=1in]{geometry}
\usepackage{setspace}
\usepackage{xcolor}
\usepackage[colorlinks=true, urlcolor=blue!60!black, linkcolor=blue!60!black]{hyperref}
\setlength{\parindent}{2em}
\setlength{\parskip}{1.25em}
\renewcommand{\baselinestretch}{1.0}
\begin{document}
\begin{center}
{\LARGE \textit{Ex parte Young}, 209 U.S. 123 (1908)\par}
\vspace{0.8em}
{\large Case Brief\par}
\end{center}
\vspace{1em}

\section{Memory Jogger}

The Eleventh Amendment does not bar a federal court from enjoining a state official from enforcing an unconstitutional state statute, because an official who acts in conflict with the supreme authority of the Federal Constitution is ``stripped of his official or representative character'' and treated as an individual rather than the State (\href{https://www.courtlistener.com/opinion/96819/ex-parte-young/}{CourtListener}; \href{https://supreme.justia.com/cases/federal/us/209/123/}{Justia}).

\section{Detailed Case Facts}

During the Progressive Era, Minnesota enacted a series of aggressive intrastate railroad rate reductions backed by severe criminal and civil penalties (\href{https://www.fjc.gov/history/cases/cases-that-shaped-the-federal-courts/ex-parte-young}{Federal Judicial Center}; \href{https://supreme.justia.com/cases/federal/us/209/123/}{Justia}). First, on September 6, 1906, Minnesota's Railroad and Warehouse Commission issued an order (effective November 15, 1906) materially lowering intrastate merchandise freight rates under the ``Western Classification,'' followed by an additional order on May 3, 1907. Under the Revised Laws of Minnesota of 1905 (\S~1987), disobeying commission orders constituted a gross misdemeanor for natural persons punishable by fines of \$2,500 to \$5,000 for a first offense and \$5,000 to \$10,000 for each subsequent offense, with matching civil forfeitures for corporations (\href{https://supreme.justia.com/cases/federal/us/209/123/}{Justia}). Second, on April 4, 1907 (effective May 1, 1907), the Minnesota legislature capped intrastate passenger fares at two cents per mile (down from three cents), making each overcharge a felony punishable by a fine up to \$5,000, imprisonment in state prison for up to five years, or both (\href{https://supreme.justia.com/cases/federal/us/209/123/}{Justia}). The railroads complied with both the 1906 commission order and the April 4 passenger act.

Third, on April 18, 1907, the legislature passed Chapter 232 (the Commodity Rate Act, effective June 1, 1907), establishing maximum intrastate carload freight rates across seven commodity classes (\href{https://supreme.justia.com/cases/federal/us/209/123/}{Justia}). Section 5 relieved the Railroad and Warehouse Commission of any statutory duty to enforce rates fixed directly by the legislature, while Section 6 made any railroad officer, director, traffic manager, agent, or employee who violated or assisted in violating the act guilty of a misdemeanor punishable by up to 90 days in county jail for each separate offense (\href{https://supreme.justia.com/cases/federal/us/209/123/}{Justia}).

On May 31, 1907---the day before Chapter 232 took effect---stockholders of nine railroads filed equity suits in the U.S. Circuit Court for the District of Minnesota, Third Division (\href{https://supreme.justia.com/cases/federal/us/209/123/}{Justia}). In the suit giving rise to the Supreme Court proceeding, Charles E. Perkins (a citizen of Iowa) and David C. Shepard (a citizen of Minnesota), shareholders of the Northern Pacific Railway Company (a Wisconsin corporation), sued the railway company, the members of the Railroad and Warehouse Commission, Minnesota Attorney General Edward T. Young, and representative shippers (\href{https://supreme.justia.com/cases/federal/us/209/123/}{Justia}). The shareholders alleged that the cumulative rate reductions were confiscatory in violation of the Fourteenth Amendment's Due Process Clause, that the draconian penalties effectively closed the courts to any test of the statutes' validity in violation of the Due Process and Equal Protection Clauses, and that the laws unlawfully burdened interstate commerce (\href{https://supreme.justia.com/cases/federal/us/209/123/}{Justia}). They sought an injunction barring Northern Pacific from implementing the rates and barring Attorney General Young from initiating enforcement actions.

\section{Procedural History}

On May 31, 1907, the U.S. Circuit Court for the District of Minnesota (Judge William Lochren) issued a temporary restraining order prohibiting Northern Pacific from publishing the April 18, 1907 commodity rates and restraining Attorney General Young from taking steps to enforce Chapter 232, while declining to disturb the already-implemented 1906 commission order and April 4 passenger act (\href{https://supreme.justia.com/cases/federal/us/209/123/}{Justia}; \href{https://www.fjc.gov/history/cases/cases-that-shaped-the-federal-courts/ex-parte-young}{Federal Judicial Center}). Young appeared specially to move to dismiss and demurred, contending that the suit against him in his official capacity was in reality an unconsented suit against the State of Minnesota barred by the Eleventh Amendment (\href{https://supreme.justia.com/cases/federal/us/209/123/}{Justia}). The Circuit Court denied the motion and overruled the demurrer.

After a hearing and presentation of evidence, the Circuit Court entered a preliminary injunction on September 23, 1907, finding the commodity rates prima facie confiscatory and enjoining both the railway from adopting Chapter 232 and Young from instituting any civil or criminal proceeding to enforce it (\href{https://supreme.justia.com/cases/federal/us/209/123/}{Justia}). The very next day, September 24, 1907, Young defied the federal injunction by petitioning a Minnesota state court (Ramsey County District Court, Second Judicial District) for an alternative writ of mandamus compelling Northern Pacific to adopt and publish the Chapter 232 rates; the state court issued the writ (\href{https://supreme.justia.com/cases/federal/us/209/123/}{Justia}).

Following an order to show cause, the federal Circuit Court adjudged Young in contempt, fined him \$100, ordered him to dismiss the state mandamus action, and committed him to the custody of the U.S. Marshal until he complied (\href{https://supreme.justia.com/cases/federal/us/209/123/}{Justia}). Young filed an original application in the U.S. Supreme Court (No. 10, Original) for writs of habeas corpus and certiorari, arguing that the Circuit Court lacked jurisdiction under the Eleventh Amendment and that its contempt order was therefore void (\href{https://supreme.justia.com/cases/federal/us/209/123/}{Justia}). The Supreme Court granted leave to file, issued a rule to show cause, heard oral argument on December 2--3, 1907, and decided the case on March 23, 1908 (\href{https://supreme.justia.com/cases/federal/us/209/123/}{Justia}).

\section{Judicial Votes}

The Supreme Court decided the case by an \textbf{8--1} vote (\href{https://supreme.justia.com/cases/federal/us/209/123/}{Justia}; \href{https://www.fjc.gov/history/cases/cases-that-shaped-the-federal-courts/ex-parte-young}{Federal Judicial Center}):

\begin{itemize}
\item \textbf{Majority (8):} Justice Rufus W. Peckham (author), joined by Chief Justice Melville W. Fuller and Justices David J. Brewer, Edward D. White, Joseph McKenna, Oliver Wendell Holmes, Jr., William R. Day, and William H. Moody.
\item \textbf{Concurrences (0):} None.
\item \textbf{Dissent (1):} Justice John Marshall Harlan.
\end{itemize}

\section{Holding}

\textbf{Judgment:} Rule to show cause discharged; petition for writs of habeas corpus and certiorari dismissed (\href{https://supreme.justia.com/cases/federal/us/209/123/}{Justia}). The U.S. Circuit Court possessed jurisdiction to enjoin Attorney General Young and to hold him in contempt for violating that injunction.

\textbf{Rule Announced:} Under the Eleventh Amendment, a federal court of equity may entertain a suit for prospective injunctive relief against a state officer clothed with some duty connected to the enforcement of an allegedly unconstitutional state statute---whether that duty arises from the challenged statute itself or from general state law---to restrain the officer from commencing civil or criminal proceedings to enforce that void enactment (\href{https://supreme.justia.com/cases/federal/us/209/123/}{Justia}). Because the State cannot authorize its officers to violate the Federal Constitution, an officer attempting to enforce an unconstitutional law is ``stripped of his official or representative character'' and cannot claim the State's sovereign immunity (\href{https://supreme.justia.com/cases/federal/us/209/123/}{Justia}). In addition, state rate-regulation penalties so enormous and severe as to intimidate regulated parties from testing the law's validity in court are unconstitutional on their face under the Fourteenth Amendment (\href{https://supreme.justia.com/cases/federal/us/209/123/}{Justia}).

\textbf{Holding vs.\ Dicta:} The Court's binding holdings are that: (1) federal-question jurisdiction existed under the Fourteenth Amendment's Due Process and Equal Protection Clauses; (2) the penalty provisions of Minnesota's rate statutes were facially unconstitutional because they effectively precluded judicial review; (3) the Eleventh Amendment did not bar a federal equity suit enjoining a state attorney general with general enforcement authority from enforcing an unconstitutional statute; and (4) the Circuit Court's contempt order was therefore valid (\href{https://supreme.justia.com/cases/federal/us/209/123/}{Justia}). By contrast, several important observations in Justice Peckham's opinion function as influential dicta (later refined in subsequent cases): his statement that a federal court may not enjoin a state court itself or a grand jury from acting; his remark that the Court did not need to decide whether the Fourteenth Amendment directly modified or limited the Eleventh Amendment; his commentary that injunctions against state rate statutes should issue only in cases ``reasonably free from doubt''; and his suggestion that juries are ill-equipped to adjudicate complex rate-confiscation questions (\href{https://supreme.justia.com/cases/federal/us/209/123/}{Justia}).

\section{Analysis of Opinions}

\subsection{Majority Opinion (Justice Peckham)}

Justice Peckham began by establishing federal-question jurisdiction, explaining that whether state-mandated railroad rates are confiscatory under the Fourteenth Amendment---and whether extreme statutory penalties deny equal protection and due process by deterring access to courts---are substantial federal questions even when they turn on complex underlying facts (\href{https://supreme.justia.com/cases/federal/us/209/123/}{Justia}). On the merits of the penalty provisions, Peckham held that subjecting railroad employees to felony imprisonment and ruinous daily fines simply for testing a rate schedule whose validity depended on technical accounting ``close[d] up all approaches to the courts,'' rendering the penalty provisions facially unconstitutional (\href{https://supreme.justia.com/cases/federal/us/209/123/}{Justia}).

Turning to the central Eleventh Amendment objection, Peckham assumed without deciding that the Fourteenth Amendment did not directly narrow the Eleventh Amendment and reaffirmed \textit{Hans v.\ Louisiana}, 134 U.S. 1 (1890), that the Eleventh Amendment bars suits against a State by its own citizens as well as citizens of other States (\href{https://supreme.justia.com/cases/federal/us/209/123/}{Justia}). Reviewing precedents from \textit{Osborn v.\ Bank of the United States}, 22 U.S. (9 Wheat.) 738 (1824), through \textit{Reagan v.\ Farmers' Loan \& Trust Co.}, 154 U.S. 362 (1894), and \textit{Smyth v.\ Ames}, 169 U.S. 466 (1898), Peckham articulated the famous ``stripping'' doctrine: when a state officer seeks to enforce an unconstitutional statute, the officer acts without valid state authority, comes into conflict with the superior authority of the Federal Constitution, and is ``stripped of his official or representative character'' (\href{https://supreme.justia.com/cases/federal/us/209/123/}{Justia}). Because the State lacks power to confer immunity from the supreme authority of the United States, enjoining the officer does not unlawfully coerce the State itself.

Peckham distinguished \textit{Fitts v.\ McGhee}, 172 U.S. 516 (1899), where state officers had no enforcement connection to the challenged bridge-toll statute. While a plaintiff may not sue a state official lacking ``some connection with the enforcement of the act'' merely to test a law's validity, Peckham held that the connection need not appear in the challenged statute itself; Attorney General Young's general statutory and common-law duties to enforce Minnesota's laws sufficed (\href{https://supreme.justia.com/cases/federal/us/209/123/}{Justia}). Nor did the injunction invade executive discretion, because an officer has no lawful discretion to enforce an unconstitutional act. Finally, Peckham rejected objections based on equity's general reluctance to enjoin criminal prosecutions and the supposed adequacy of legal defenses in state court: once a federal equity court first acquires jurisdiction over a constitutional challenge, it may enjoin a party from commencing subsequent state civil or criminal enforcement suits (though it may not enjoin state courts or grand juries themselves), and awaiting criminal prosecution was an inadequate remedy at law given the crushing penalties and technical complexity of rate litigation (\href{https://supreme.justia.com/cases/federal/us/209/123/}{Justia}).

\subsection{Dissenting Opinion (Justice Harlan)}

In a forceful lone dissent, Justice Harlan argued that the majority's distinction between the State and its Attorney General was an unsustainable fiction that threatened to ``practically obliterate the Eleventh Amendment'' (\href{https://supreme.justia.com/cases/federal/us/209/123/}{Justia}). Because ``the intangible thing called a State \ldots\ can never appear or be represented or known in any court in a litigated case except by and through its officers,'' suing Young ``as, and only because he was, Attorney General of Minnesota'' to stop him from filing a mandamus suit in the State's own courts was, in legal and practical effect, a suit against the State of Minnesota (\href{https://supreme.justia.com/cases/federal/us/209/123/}{Justia}).

Harlan contended that the majority departed from \textit{In re Ayers}, 123 U.S. 443 (1887), and ``frittered away'' \textit{Fitts v.\ McGhee}, 172 U.S. 516 (1899), both of which recognized a sharp line between enjoining an officer from committing a personal trespass upon property and forbidding a state's chief law officer from invoking the jurisdiction of the State's own courts (\href{https://supreme.justia.com/cases/federal/us/209/123/}{Justia}). He warned that empowering subordinate federal courts to ``put manacles upon the State'' would reduce States to the status of ``dependencies'' or provinces (\href{https://supreme.justia.com/cases/federal/us/209/123/}{Justia}). Stressing that state courts share an equal constitutional oath and obligation to enforce federal rights under the Supremacy Clause, Harlan maintained that Northern Pacific should have raised its Fourteenth Amendment defenses in Minnesota's courts, with ultimate review by writ of error in the U.S. Supreme Court (\href{https://supreme.justia.com/cases/federal/us/209/123/}{Justia}).

\section{Examples: Future Applications}

\begin{enumerate}
\item \textbf{Same-Side (Permitted under \textit{Ex parte Young}):} A state legislature enacts a statute banning private clinics from providing a federally protected medical procedure and directs the state health commissioner to revoke the licenses of violating clinics and assess \$50,000 daily fines. A clinic sues the health commissioner in federal court seeking a prospective injunction barring enforcement. \textit{Why it falls on the same side:} The suit names a state official who has a direct enforcement connection to the challenged statute and seeks purely prospective equitable relief to stop an ongoing or threatened violation of federal law, fitting squarely within \textit{Ex parte Young} (\href{https://supreme.justia.com/cases/federal/us/209/123/}{Justia}).

\item \textbf{Same-Side (Permitted under \textit{Ex parte Young}):} A state environmental agency director refuses to comply with mandatory federal Clean Air Act permitting standards, imposing ongoing illegal restrictions on a utility. The utility sues the director in federal court for declaratory and prospective injunctive relief requiring future compliance, even though compliance will require the state agency to spend money on updated monitoring equipment. \textit{Why it falls on the same side:} Under \textit{Ex parte Young} and \href{https://supreme.justia.com/cases/federal/us/415/651/}{\textit{Edelman v.\ Jordan}}, 415 U.S. 651 (1974), prospective injunctions ordering state officers to conform future conduct to federal law are permissible even when they have an ``ancillary effect'' on the state treasury, so long as no detailed federal statutory remedial scheme displaces \textit{Young}.

\item \textbf{Opposite-Side (Barred by the Eleventh Amendment):} State welfare recipients sue a state human services director in federal court seeking an order compelling the director to pay two years of wrongfully withheld past welfare benefits out of the state treasury. \textit{Why it falls on the opposite side:} As the Supreme Court later held in \href{https://supreme.justia.com/cases/federal/us/415/651/}{\textit{Edelman v.\ Jordan}}, 415 U.S. 651 (1974), \textit{Ex parte Young} permits only prospective injunctive relief; an award of retroactive monetary liability payable from public funds in the state treasury is indistinguishable from a damages judgment against the State itself and is barred by the Eleventh Amendment.

\item \textbf{Opposite-Side (Barred by the Eleventh Amendment):} Residents of a state-run psychiatric facility sue state officials in federal court seeking a prospective injunction compelling the officials to comply with a \emph{state} mental-health statute that provides stronger protections than federal law. \textit{Why it falls on the opposite side:} As the Court held in \href{https://www.law.cornell.edu/supremecourt/text/465/89}{\textit{Pennhurst State School \& Hospital v.\ Halderman}}, 465 U.S. 89 (1984), the \textit{Ex parte Young} fiction rests on the need to vindicate the supremacy of \emph{federal} law; that rationale disappears when a plaintiff asks a federal court to instruct state officials on how to conform their conduct to state law.

\item \textbf{Fence-Sitter (Borderline / Contested Application):} A state passes a statute restricting a constitutional right but expressly forbids the state attorney general and all executive officials from enforcing it, delegating enforcement exclusively to private civil lawsuits in state court while retaining ancillary state licensing-board oversight over regulated professionals. Regulated providers sue state-court clerks, judges, the attorney general, and state licensing officials in federal court under \textit{Ex parte Young}. \textit{Why it is a fence-sitter:} Modeled on \href{https://supreme.justia.com/cases/federal/us/595/21-463/}{\textit{Whole Woman's Health v.\ Jackson}}, 595 U.S. 30 (2021), this scenario straddles the exact boundary of \textit{Ex parte Young}: suits against state judges, clerks, and an attorney general stripped of enforcement authority fail under \textit{Young}'s ban on enjoining state courts and its ``some connection with enforcement'' rule, whereas claims against executive licensing officials who may indirectly enforce the law through disciplinary sanctions can survive---illustrating how easily states can test the formalist limits of \textit{Young}'s officer-suit fiction.
\end{enumerate}

\section{Critique}

\textit{Ex parte Young} is one of the most celebrated yet conceptually paradoxical decisions in American constitutional law. Its central analytical vulnerability is the famous ``double fiction'' (or paradox) embedded in Justice Peckham's opinion: for purposes of the Fourteenth Amendment, Attorney General Young's threatened enforcement constitutes ``state action'' triggering constitutional limits, yet for purposes of the Eleventh Amendment, the exact same official act is treated as an unauthorized private wrong stripped of state character (\href{https://supreme.justia.com/cases/federal/us/209/123/}{Justia}; \href{https://www.fjc.gov/history/cases/cases-that-shaped-the-federal-courts/ex-parte-young}{Federal Judicial Center}).

From an \textbf{originalist and textualist perspective}, scholars and jurists have divided sharply. Some formalist critics echo Justice Harlan's dissent, arguing that \textit{Young} elevates pleading form over constitutional substance by allowing federal courts to control state policy through injunctions against officers, thereby eroding the structural sovereign immunity recognized in \textit{Hans v.\ Louisiana} and \href{https://supreme.justia.com/cases/federal/us/527/706/}{\textit{Alden v.\ Maine}}, 527 U.S. 706 (1999). Recent historical scholarship by Judge Andrew S. Oldham, Alexander Steene, and James Tienken similarly contends that \textit{Young}'s non-statutory equitable cause of action would not have been recognized in the English Court of Chancery in 1789 and sits uneasily alongside modern textualist skepticism of judicially implied causes of action (\href{https://scholarlycommons.law.northwestern.edu/cgi/viewcontent.cgi?article=1648&context=nulr}{Northwestern University Law Review}). Conversely, other prominent originalist and historical scholars defend \textit{Young} without relying on a fiction: John Harrison argues that \textit{Young} was simply a traditional equitable anti-suit injunction allowing a prospective defendant to raise a federal defense affirmatively without offending sovereign immunity (\href{https://review.law.stanford.edu/print/article/ex-parte-young/}{Stanford Law Review}), while James Pfander and Jacob Wentzel trace officer-suit remedies to deep common-law and equitable traditions (\href{https://review.law.stanford.edu/wp-content/uploads/sites/3/2020/05/Pfander-Wentzel-72-Stan.-L.-Rev.-1269.pdf}{Stanford Law Review}).

From a \textbf{progressive and doctrinal perspective}, \textit{Young}'s reputation has undergone a dramatic historical reversal. Decided only three years after \textit{Lochner v.\ New York} (1905), \textit{Young} was initially denounced by Progressives as an instrument of ``government by injunction'' that empowered federal judges to paralyze state economic and railroad regulation---prompting Congress in 1910 to require three-judge district courts for injunctions against state statutes (\href{https://www.fjc.gov/history/cases/cases-that-shaped-the-federal-courts/ex-parte-young}{Federal Judicial Center}). Yet by the mid-twentieth century, as legal historian Edward A. Purcell, Jr.\ has observed, \textit{Young} became the indispensable procedural backbone of the civil rights movement, enabling federal courts to enjoin state-mandated racial segregation, malapportionment, and constitutional abuses (\href{https://www.fjc.gov/history/cases/cases-that-shaped-the-federal-courts/ex-parte-young}{Federal Judicial Center}). Modern doctrinal critics, however, warn that because \textit{Young} rests on an officer-suit formalism rather than a direct constitutional right to sue States, the Supreme Court has repeatedly narrowed it---precluding \textit{Young} suits where Congress created a detailed statutory remedial scheme (\href{https://supreme.justia.com/cases/federal/us/517/44/}{\textit{Seminole Tribe of Florida v.\ Florida}}, 517 U.S. 44 (1996)), barring retroactive relief (\href{https://supreme.justia.com/cases/federal/us/415/651/}{\textit{Edelman v.\ Jordan}}, 415 U.S. 651 (1974)), and leaving gaps when States outsource enforcement to private litigants (\href{https://supreme.justia.com/cases/federal/us/595/21-463/}{\textit{Whole Woman's Health v.\ Jackson}}, 595 U.S. 30 (2021)).

At bottom, \textit{Ex parte Young} reflects an enduring structural compromise between two competing constitutional values: \textbf{federal constitutional supremacy and effective judicial remedies for individual rights} versus \textbf{state sovereign dignity, fiscal autonomy, and comity toward state courts} (\href{https://www.fjc.gov/history/cases/cases-that-shaped-the-federal-courts/ex-parte-young}{Federal Judicial Center}; \href{https://scholarship.shu.edu/cgi/viewcontent.cgi?article=1034&context=circuit_review}{Seton Hall Circuit Review}).

\section{Key Quotations}

\begin{quote}
``If the act which the state Attorney General seeks to enforce be a violation of the Federal Constitution, the officer, in proceeding under such enactment, comes into conflict with the superior authority of that Constitution, and he is, in that case, stripped of his official or representative character, and is subjected in his person to the consequences of his individual conduct. The State has no power to impart to him any immunity from responsibility to the supreme authority of the United States.''
\end{quote}
\noindent --- Justice Rufus W. Peckham, Majority Opinion, 209 U.S.\ at 159--60 (\href{https://www.courtlistener.com/opinion/96819/ex-parte-young/}{CourtListener}; \href{https://supreme.justia.com/cases/federal/us/209/123/}{Justia}).

\begin{quote}
``In making an officer of the State a party defendant in a suit to enjoin the enforcement of an act alleged to be unconstitutional, it is plain that such officer must have some connection with the enforcement of the act, or else it is merely making him a party as a representative of the State, and thereby attempting to make the State a party.''
\end{quote}
\noindent --- Justice Rufus W. Peckham, Majority Opinion, 209 U.S.\ at 157 (\href{https://www.courtlistener.com/opinion/96819/ex-parte-young/}{CourtListener}; \href{https://supreme.justia.com/cases/federal/us/209/123/}{Justia}).

\begin{quote}
``We hold, therefore, that the provisions of the acts relating to the enforcement of the rates, either for freight or passengers, by imposing such enormous fines and possible imprisonment as a result of an unsuccessful effort to test the validity of the laws themselves, are unconstitutional on their face, without regard to the question of the insufficiency of those rates.''
\end{quote}
\noindent --- Justice Rufus W. Peckham, Majority Opinion, 209 U.S.\ at 148 (\href{https://www.courtlistener.com/opinion/96819/ex-parte-young/}{CourtListener}; \href{https://supreme.justia.com/cases/federal/us/209/123/}{Justia}).

\begin{quote}
``It would therefore seem clear that, within the true meaning of the Eleventh Amendment, the suit brought in the Federal court was one, in legal effect, against the State -- as much so as if the State had been formally named on the record as a party -- and therefore it was a suit to which, under the Amendment, so far as the State or its Attorney General was concerned, the judicial power of the United States did not and could not extend.''
\end{quote}
\noindent --- Justice John Marshall Harlan, Dissenting Opinion, 209 U.S.\ at 174 (\href{https://www.courtlistener.com/opinion/96819/ex-parte-young/}{CourtListener}; \href{https://supreme.justia.com/cases/federal/us/209/123/}{Justia}).

\begin{quote}
``This principle, if firmly established, would work a radical change in our governmental system. It would inaugurate a new era in the American judicial system and in the relations of the national and state governments. It would enable the subordinate Federal courts to supervise and control the official action of the States as if they were `dependencies' or provinces.''
\end{quote}
\noindent --- Justice John Marshall Harlan, Dissenting Opinion, 209 U.S.\ at 175 (\href{https://www.courtlistener.com/opinion/96819/ex-parte-young/}{CourtListener}; \href{https://supreme.justia.com/cases/federal/us/209/123/}{Justia}).


\end{document}
